\documentclass[12pt,a4paper]{article}
\usepackage[margin=25mm, headheight=15pt, headsep=10pt]{geometry}
\usepackage{fontspec}
\usepackage[table]{xcolor} % Note: table option is required for rowcolors
\usepackage{titlesec}
\usepackage{tocloft}
\usepackage{fancyhdr}
\usepackage{booktabs}
\usepackage{tcolorbox}
\tcbuselibrary{skins, breakable}
\usepackage[colorlinks=true, linkcolor=emerald, urlcolor=emerald, bookmarks=true]{hyperref}

\definecolor{emerald}{RGB}{0,102,51}
\definecolor{darkred}{RGB}{139,0,0}
\definecolor{lightgray}{RGB}{245,245,245}

\usepackage[nil]{babel}
\babelprovide[import, main]{bengali}
\babelprovide[import]{arabic}
\babelfont{rm}[Renderer=HarfBuzz,Scale=1.0]{Noto Serif Bengali}
\babelfont{sf}[Renderer=HarfBuzz]{Noto Sans Bengali}
\babelfont{tt}[Renderer=HarfBuzz]{Noto Sans Bengali}
\babelfont[arabic]{rm}[Renderer=HarfBuzz,Scale=1.1]{Noto Naskh Arabic}

% Section styling
\titleformat{\section}{\Large\bfseries\color{emerald}}{\thesection.}{0.5em}{}
\titleformat{\subsection}{\large\bfseries\color{emerald!85!black}}{\thesubsection.}{0.5em}{}
\titleformat{\subsubsection}{\normalsize\bfseries\color{emerald!70!black}}{\thesubsubsection.}{0.5em}{}
\titlespacing*{\section}{0pt}{18pt}{8pt}

% Fancy Header/Footer setup
\pagestyle{fancy}
\fancyhf{} % clear all header and footer fields
\fancyhead[L]{\color{emerald}\small\textbf{\nouppercase{\leftmark}}}
\fancyfoot[L]{\scriptsize\color{black!50}{\lat{AI}}-সংকলিত, আলেম-যাচাইকৃত নয়}
\fancyfoot[C]{\color{emerald!70!black}\thepage}
\renewcommand{\headrulewidth}{0.4pt}
\renewcommand{\headrule}{\hbox to\headwidth{\color{emerald}\leaders\hrule height \headrulewidth\hfill}}
\renewcommand{\footrulewidth}{0pt}

% Custom boxes
% 1. Ayat/Hadith Box
\newtcolorbox{ayatbox}[1][]{
    enhanced, breakable, 
    colback=emerald!5, 
    colframe=emerald, 
    boxrule=0.5pt, 
    arc=3pt, 
    left=10pt, right=10pt, top=10pt, bottom=10pt,
    #1
}

% 2. Quote/Translation Box
\newtcolorbox{quotebox}[1][]{
    enhanced, breakable,
    frame hidden,
    colback=lightgray,
    borderline west={3pt}{0pt}{emerald},
    left=15pt, right=10pt, top=10pt, bottom=10pt,
    sharp corners,
    #1
}

% 3. AI-generated notice box
\newtcolorbox{warnbox}[1][]{
    enhanced, breakable,
    colback=red!4,
    colframe=darkred,
    boxrule=0.8pt,
    arc=3pt,
    left=10pt, right=10pt, top=10pt, bottom=10pt,
    #1
}

% Commands
\newcommand{\arab}[1]{\foreignlanguage{arabic}{#1}}
\newcommand{\kib}[1]{\textcolor{emerald}{\textbf{#1}}}

% Latin (English) fallback: Noto Serif Bengali has NO Latin glyphs
\newfontfamily\latfont[Renderer=HarfBuzz]{Noto Serif}
\newcommand{\lat}[1]{{\latfont #1}}

\setlength{\parskip}{5pt}
\setlength{\parindent}{0pt}

% Fix for missing bullet in Noto Serif Bengali
\renewcommand{\labelitemi}{$\bullet$}



\begin{document}
\begin{center}{\Large\bfseries\color{emerald} ক্ষমা করা পাপ — কিয়ামতের দিন কী হয়? (হাদিস ও আসার)}\end{center}
\vspace{1em}
\section*{ক্ষমা করা পাপ — কিয়ামতের দিন কী হয়? (হাদিস ও আসার)}

\begin{warnbox}
 \textbf{গুরুত্বপূর্ণ নোটিশ (\lat{Important} \lat{Notice}):} এই কন্টেন্টটি \lat{AI} (কৃত্রিম বুদ্ধিমত্তা) দিয়ে প্রস্তুত ও সংকলিত। \lat{AI} কঠোর সূত্র-মান নীতি মেনে শুধু নির্ভরযোগ্য উৎস থেকে তথ্য সংগ্রহ করেছে — কেবল সহীহ/হাসান হাদিস ও সহীহ সনদের আসার রাখা হয়েছে, দুর্বল সনদ বাদ দেওয়া হয়েছে। তবে এটি কোনো আলেম বা মুহাদ্দিস কর্তৃক যাচাইকৃত নয়।
\end{warnbox}

\begin{quote}
\textbf{পড়ার নিয়ম:} প্রতিটি হাদিসে — সূত্র কার্ড (সনদের মান, গ্রন্থ, অধ্যায়, নম্বর) $\rightarrow$ পূর্ণ অনুবাদ $\rightarrow$ শব্দের ব্যাখ্যা $\rightarrow$ সংশ্লিষ্ট দলিল। শেষে প্রচলিত-কিন্তু-দুর্বল কনটেন্টের সতর্কতা টেবিল। \\
\textbf{মূলনীতি:} শুধু সহীহ/হাসান; দুর্বল সনদ লেবেল দিয়ে নয়, সম্পূর্ণ বাদ।
\end{quote}

\medskip\hrule\medskip

\subsection*{হাদিস ১ — নাজওয়া হাদিস: মুমিনের পাপ একান্তে স্মরণ করিয়ে ক্ষমা (সবচেয়ে গুরুত্বপূর্ণ)}

\subsubsection*{সূত্র কার্ড}

\begin{table}[h]\centering\small
\begin{tabular}{ll}
বিষয় & বিবরণ \\
\hline
\textbf{বর্ণনাকারী} & ইবনে উমার (রাঃ) — সাফওয়ান ইবনে মুহরিজ (রহ.)-এর কাছে \\
\textbf{গ্রন্থ} & সহীহ বুখারি ২৪৪১, ৪৬৮৫, ৬০৭০, ৭৫১৪; সহীহ মুসলিম ২৭৬৮ \\
\textbf{অধ্যায়} & বুখারি: কিতাবুল মাযালিম — \lat{“}যালিমদের উপর আল্লাহর লানত\lat{”}; মুসলিম: কিতাবুত তাওবাহ — \lat{“}হত্যাকারীর তাওবা কবুল, হত্যা যত বেশি হোক\lat{”} \\
\textbf{সনদের মান} & \textbf{সহীহ} (বুখারি-মুসলিম) — সর্বোচ্চ মান \\
\hline
\end{tabular}
\end{table}

\subsubsection*{পূর্ণ অনুবাদ}

সাফওয়ান ইবনে মুহরিজ (রহ.) বলেন: একদা আমি ইবনে উমার (রাঃ)-এর সাথে হাঁটছিলাম, তাঁর হাত ধরে। তখন এক ব্যক্তি এসে জিজ্ঞেস করল: \lat{“}নাজওয়া (আল্লাহর সাথে বান্দার একান্ত গোপন কথোপকথন) সম্পর্কে আপনি রাসূলুল্লাহ (সা.)-এর কাছ থেকে কী শুনেছেন?\lat{”} ইবনে উমার (রাঃ) বললেন: আমি রাসূলুল্লাহ (সা.)-কে বলতে শুনেছি:

\textbf{\lat{“}আল্লাহ কিয়ামতের দিন মুমিনকে তাঁর (নিজের) নিকটে টেনে নেবেন, তার উপর তাঁর আড়াল (কানাফ) রেখে তাকে ঢেকে দেবেন এবং বলবেন: "তুমি কি এই পাপটি চেনো? তুমি কি এই পাপটি চেনো?" সে বলবে: "হ্যাঁ, হে আমার রব!" এভাবে আল্লাহ তাকে তার সব পাপ স্বীকার করাবেন। সে মনে করবে সে তো ধ্বংস হয়ে গেল। তখন আল্লাহ বলবেন: "আমি দুনিয়ায় এগুলো তোমার উপর গোপন রেখেছিলাম এবং আজ আমি তোমার জন্য এগুলো ক্ষমা করে দিলাম।" এরপর তাকে তার নেকির আমলনামা দেওয়া হবে। আর কাফির ও মুনাফিকদের ব্যাপারে ঘোষণা করা হবে — (ফেরেশতা ও সাক্ষীরা বলবে): "এরা তারাই, যারা তাদের রবের বিরুদ্ধে মিথ্যা বলেছিল। জেনে রাখো, যালিমদের উপর আল্লাহর লানত!"\lat{”}}

(সহীহ বুখারি ২৪৪১; সহীহ মুসলিম ২৭৬৮)

\subsubsection*{শব্দের ব্যাখ্যা}

\begin{table}[h]\centering\small
\begin{tabular}{ll}
শব্দ & অর্থ \\
\hline
\textbf{আন-নাজওয়া (\arab{النجوى})} & গোপন কথোপকথন — আল্লাহ ও বান্দার একান্ত আলাপ \\
\textbf{ইউদনিল মুমিন} & মুমিনকে কাছে টেনে নেন \\
\textbf{কানাফুহু (\arab{كنفه})} & তাঁর আড়াল/পর্দা — যা ঢেকে রাখে \\
\textbf{ইউক্বাররিরুহু} & স্বীকার করান — পাপগুলো স্বীকার করান \\
\textbf{সাতারতুহা 'আলাইকা} & আমি তোমার উপর (এগুলো) গোপন রেখেছিলাম \\
\textbf{কিতাবা হাসানাতিহি} & তার নেকির আমলনামা \\
\hline
\end{tabular}
\end{table}

\subsubsection*{সংশ্লিষ্ট দলিল}

\begin{itemize}
\item \textbf{বুখারি ৪৬৮৫, ৬০৭০, ৭৫১৪:} একই হাদিসের আরও বর্ণনা — সব সহীহ।
\item \textbf{ইবনে হাজারের (রহ.) ব্যাখ্যা (ফাতহুল বারি):} এই হাদিস প্রমাণ করে — \textbf{আল্লাহ কিয়ামতের দিন মুমিনদের পাপ গোপন রাখবেন}; কাফির ও মুনাফিকদের ক্ষেত্রে এই গোপন রাখার বিধান নেই — তাদের পাপ প্রকাশ করা হবে।
\item \textbf{গুরুত্বপূর্ণ পয়েন্ট:} পাপ \lat{“}স্মরণ করানো\lat{”} হচ্ছে — কিন্তু শাস্তি বা অপমানের জন্য নয়; বরং বান্দা নিজে স্বীকার করুক ও আল্লাহর অনুগ্রহ বুঝুক — তারপর ক্ষমা। এখানেই ক্ষমা করা পাপের কিয়ামতের অবস্থান স্পষ্ট: \textbf{মুমিনের ক্ষমা করা পাপ আর কোনো বোঝা নয় — তাকে নেকির আমলনামাই দেওয়া হয়।}
\end{itemize}
\medskip\hrule\medskip

\subsection*{হাদিস ২ — সহজ হিসাবের ব্যাখ্যা: \lat{“}যার হিসাব পুঙ্খানুপুঙ্খ নেওয়া হবে, সে শাস্তি পাবে\lat{”}}

\subsubsection*{সূত্র কার্ড}

\begin{table}[h]\centering\small
\begin{tabular}{ll}
বিষয় & বিবরণ \\
\hline
\textbf{বর্ণনাকারী} & আয়েশা (রাঃ) \\
\textbf{গ্রন্থ} & সহীহ বুখারি ৬৫৩৬; সহীহ মুসলিম ২৮৭৬ \\
\textbf{অধ্যায়} & বুখারি: কিতাবুর রিকাক — \lat{“}যার হিসাব পুঙ্খানুপুঙ্খ নেওয়া হবে সে শাস্তি পাবে\lat{”}; মুসলিম: কিতাবুল জান্নাহ — \lat{“}হিসাব প্রমাণ\lat{”} \\
\textbf{সনদের মান} & \textbf{সহীহ} (বুখারি-মুসলিম) \\
\hline
\end{tabular}
\end{table}

\subsubsection*{পূর্ণ অনুবাদ}

আয়েশা (রাঃ) বলেন: রাসূলুল্লাহ (সা.) বলেছেন: \textbf{\lat{“}কিয়ামতের দিন যার হিসাব পুঙ্খানুপুঙ্খভাবে নেওয়া হবে, সে শাস্তি পাবে।\lat{”}} আমি বললাম: আল্লাহ কি বলেননি — \lat{“}সে তো সহজ হিসাবে হিসাব পেশ করবে\lat{”} (৮৪:৮)? তিনি বললেন: \textbf{\lat{“}সেটি তো আরয (আমলের পেশকরণ) — পুঙ্খানুপুঙ্খ হিসাব নয়। যার হিসাব পুঙ্খানুপুঙ্খভাবে নেওয়া হবে, সে শাস্তি পাবে।\lat{”}}

(সহীহ বুখারি ৬৫৩৬; সহীহ মুসলিম ২৮৭৬)

\subsubsection*{শব্দের ব্যাখ্যা}

\begin{table}[h]\centering\small
\begin{tabular}{ll}
শব্দ & অর্থ \\
\hline
\textbf{নুকিশা (\arab{نوقش})} & পুঙ্খানুপুঙ্খভাবে জিজ্ঞাসাবাদ করা — প্রতিটি বিষয়ে খুঁটিনাটি \\
\textbf{উয্যিবা (\arab{عذّب})} & শাস্তি পাবে — ধ্বংস হবে \\
\textbf{আল-আরয (\arab{العرض})} & পেশকরণ — আমল দেখানো, খুঁটিনাটি জিজ্ঞাসা নয় \\
\hline
\end{tabular}
\end{table}

\subsubsection*{সংশ্লিষ্ট দলিল}

\begin{itemize}
\item \textbf{আয়াত (৮৪:৮):} \lat{“}সে তো সহজ হিসাবে হিসাব পেশ করবে\lat{”} — নবীজির (সা.) ব্যাখ্যা: এটি আরয।
\item \textbf{হাদিস (বুখারি ২৪৪১; মুসলিম ২৭৬৮):} নাজওয়া — মুমিনের পাপ একান্তে স্মরণ করানো হয়, তারপর ক্ষমা ও নেকির আমলনামা।
\item \textbf{নববি (রহ.) (শরহে মুসলিম):} মুমিনের হিসাব সহজ হবে — ক্ষমা করা পাপ পুনরায় হিসাবের বিষয় নয়।
\end{itemize}
\medskip\hrule\medskip

\subsection*{হাদিস ৩ — যে মুসলিমের দোষ গোপন রাখে, আল্লাহ তাকে দুনিয়া ও আখিরাতে গোপন রাখবেন}

\subsubsection*{সূত্র কার্ড}

\begin{table}[h]\centering\small
\begin{tabular}{ll}
বিষয় & বিবরণ \\
\hline
\textbf{বর্ণনাকারী} & আবু হুরাইরা (রাঃ) \\
\textbf{গ্রন্থ} & সহীহ মুসলিম ২৬৯৯ \\
\textbf{অধ্যায়} & কিতাবুজ যিকরি ওদ দোয়া — \lat{“}কুরআন তিলাওয়াত ও যিকরের মজলিসের ফজিলত\lat{”} \\
\textbf{সনদের মান} & \textbf{সহীহ} \\
\hline
\end{tabular}
\end{table}

\subsubsection*{পূর্ণ অনুবাদ}

আবু হুরাইরা (রাঃ) থেকে বর্ণিত — রাসূলুল্লাহ (সা.) বলেছেন: \textbf{\lat{“}যে ব্যক্তি কোনো মুসলিমের দোষ গোপন রাখে, আল্লাহ তার দোষ দুনিয়া ও আখিরাতে গোপন রাখবেন। আল্লাহ বান্দার সাহায্যে থাকেন, যতক্ষণ বান্দা তার ভাইয়ের সাহায্যে থাকে...\lat{”}}

(সহীহ মুসলিম ২৬৯৯)

\subsubsection*{সংশ্লিষ্ট দলিল}

\begin{itemize}
\item \textbf{হাদিস (বুখারি ৬০৬৯):} \lat{“}আমার সব উম্মত ক্ষমা পাবে, কিন্তু মুজাহির (যারা পাপ প্রকাশ করে) ছাড়া।\lat{”}
\item \textbf{হাদিস (বুখারি ২৪৪১; মুসলিম ২৭৬৮):} আল্লাহ দুনিয়ায় পাপ গোপন রাখলে আখিরাতেও গোপন রাখেন — \lat{“}আমি দুনিয়ায় এগুলো তোমার উপর গোপন রেখেছিলাম এবং আজ ক্ষমা করে দিলাম।\lat{”}
\end{itemize}
\medskip\hrule\medskip

\subsection*{হাদিস ৪ — মুজাহির ছাড়া সব উম্মত ক্ষমা পাবে}

\subsubsection*{সূত্র কার্ড}

\begin{table}[h]\centering\small
\begin{tabular}{ll}
বিষয় & বিবরণ \\
\hline
\textbf{বর্ণনাকারী} & আবু হুরাইরা (রাঃ) \\
\textbf{গ্রন্থ} & সহীহ বুখারি ৬০৬৯ \\
\textbf{অধ্যায়} & কিতাবুল আদাব — \lat{“}মুমিন নিজের পাপ নিজে গোপন রাখবে\lat{”} \\
\textbf{সনদের মান} & \textbf{সহীহ} \\
\hline
\end{tabular}
\end{table}

\subsubsection*{পূর্ণ অনুবাদ}

আবু হুরাইরা (রাঃ) থেকে বর্ণিত — রাসূলুল্লাহ (সা.) বলেছেন: \textbf{\lat{“}আমার সব উম্মত ক্ষমা পাবে, কিন্তু মুজাহির (যারা পাপ প্রকাশ করে) ছাড়া। পাপ প্রকাশের উদাহরণ — কেউ রাতে এমন পাপ করল, আল্লাহ তা ঢেকে রাখলেন; কিন্তু সকালে সে বলল: 'ওহে অমুক! গত রাতে আমি এমন এমন করেছি।' সারারাত তার রব তাকে ঢেকে রেখেছিলেন, আর সকালে সে নিজেই আল্লাহর পর্দা ছিঁড়ে ফেলল।\lat{”}}

(সহীহ বুখারি ৬০৬৯)

\subsubsection*{সংশ্লিষ্ট দলিল}

\begin{itemize}
\item \textbf{হাদিস (মুসলিম ২৬৯৯):} দোষ গোপন রাখার পুরস্কার — দুনিয়া ও আখিরাতে গোপন।
\item \textbf{ব্যাখ্যা (ইবনে হাজার):} নিজের পাপ প্রকাশ করা আল্লাহর গোপন রাখার অনুগ্রহকে বরবাদ করে; পাপ গোপন রাখা — তাওবা ও আখিরাতের গোপন রাখার সাথে সম্পর্কিত।
\end{itemize}
\medskip\hrule\medskip

\subsection*{হাদিস ৫ — মুফলিস (দেউলিয়া) হাদিস: মানুষের হক ক্ষমার সাধারণ বিধানে আসে না}

\subsubsection*{সূত্র কার্ড}

\begin{table}[h]\centering\small
\begin{tabular}{ll}
বিষয় & বিবরণ \\
\hline
\textbf{বর্ণনাকারী} & আবু হুরাইরা (রাঃ) \\
\textbf{গ্রন্থ} & সহীহ মুসলিম ২৫৮১ \\
\textbf{অধ্যায়} & কিতাবুল বিরর — \lat{“}জুলুম নিষিদ্ধ\lat{”} \\
\textbf{সনদের মান} & \textbf{সহীহ} \\
\hline
\end{tabular}
\end{table}

\subsubsection*{পূর্ণ অনুবাদ}

আবু হুরাইরা (রাঃ) থেকে বর্ণিত — রাসূলুল্লাহ (সা.) জিজ্ঞেস করলেন: \textbf{\lat{“}তোমরা কি জানো, মুফলিস (দেউলিয়া/কপর্দকহীন) কে?\lat{”}} সাহাবিরা বললেন: আমাদের মধ্যে মুফলিস সে, যার কাছে দিরহামও নেই, সম্পদও নেই। তিনি বললেন: \textbf{\lat{“}আমার উম্মতের মুফলিস সে, যে কিয়ামতের দিন নামাজ, রোজা ও যাকাত নিয়ে আসবে, কিন্তু সে এদিকে ওই ব্যক্তিকে গালি দিয়েছে, ওই ব্যক্তির উপর অপবাদ দিয়েছে, ওই ব্যক্তির সম্পদ খেয়েছে, ওই ব্যক্তির রক্তপাত করেছে, ওই ব্যক্তিকে মারধর করেছে — তখন তার নেকিগুলো থেকে এদের প্রত্যেককে দেওয়া হবে। আর তার নেকি ফুরিয়ে গেলে, তার উপর (যাদের হক নষ্ট করেছে) তাদের পাপ চাপিয়ে দেওয়া হবে; এরপর তাকে জাহান্নামে নিক্ষেপ করা হবে।\lat{”}}

(সহীহ মুসলিম ২৫৮১)

\subsubsection*{সংশ্লিষ্ট দলিল}

\begin{itemize}
\item \textbf{হাদিস (বুখারি ২৪৪৯):} \lat{“}যার কাছে কারো জুলুমের হক আছে, সে যেন আজই তা ফেরত দেয়/ক্ষমা আদায় করে — আগামীকাল এমন দিন আসবে যেদিন দিরহাম-দীনার থাকবে না; তার নেকি থেকে নেওয়া হবে।\lat{”}
\item \textbf{হাদিস (বুখারি ৩১৮৮):} \lat{“}প্রত্যেক বিশ্বাসঘাতকের জন্য কিয়ামতের দিন একটি পতাকা উঁচু করা হবে — যাতে তার বিশ্বাসঘাতকতা চেনা যায়।\lat{”}
\item \textbf{গুরুত্বপূর্ণ পার্থক্য:} \textbf{আল্লাহর হক} — তাওবা, ইস্তিগফার, নেক আমলে ক্ষমা হয়; \textbf{মানুষের হক} — ফেরত দেওয়া/ক্ষমা আদায় করা ছাড়া ক্ষমা হয় না। এই দুইয়ের পার্থক্য বুঝলে \lat{“}ক্ষমা করা পাপের কিয়ামতের হিসাব\lat{”} প্রশ্নটি পরিষ্কার হয়।
\end{itemize}
\medskip\hrule\medskip

\subsection*{হাদিস ৬ — ইসলাম/হিজরত/হজ: যা আগে ছিল তা মুছে দেয় (নতুন মুসলিমের ক্ষেত্রে)}

\subsubsection*{সূত্র কার্ড}

\begin{table}[h]\centering\small
\begin{tabular}{ll}
বিষয় & বিবরণ \\
\hline
\textbf{বর্ণনাকারী} & আমর ইবনে আস (রাঃ) \\
\textbf{গ্রন্থ} & সহীহ মুসলিম ১২১ \\
\textbf{অধ্যায়} & কিতাবুল ঈমান — \lat{“}ইসলাম যা আগে ছিল তা ধ্বংস করে দেয়, হিজরত ও হজও\lat{”} \\
\textbf{সনদের মান} & \textbf{সহীহ} \\
\hline
\end{tabular}
\end{table}

\subsubsection*{পূর্ণ অনুবাদ}

আমর ইবনে আস (রাঃ) মৃত্যুশয্যায় কাঁদছিলেন। তাঁর ছেলে বলল: \lat{“}রাসূলুল্লাহ (সা.) কি আপনাকে সুসংবাদ দেননি?\lat{”} তিনি বললেন: \lat{“}হ্যাঁ — সাক্ষ্য যে আল্লাহ ছাড়া কোনো ইলাহ নেই এবং মুহাম্মাদ আল্লাহর রাসূল। আমি তিনটি অবস্থা পার করেছি... আল্লাহ আমার অন্তরে ইসলামের ভালোবাসা দিলেন, আমি নবীজির (সা.) কাছে এসে বললাম: আপনার ডান হাত বাড়ান, আমি আপনার কাছে বাইয়াত করি। আমি হাত গুটিয়ে নিলাম। নবীজি (সা.) বললেন: কী হয়েছে? আমি বললাম: আমি শর্ত রাখতে চাই। তিনি বললেন: কী শর্ত? আমি বললাম: যেন আমি ক্ষমা পাই। তিনি বললেন: \textbf{তুমি কি জানো না যে, ইসলাম যা আগে ছিল তা মুছে দেয়? হিজরত যা আগে ছিল তা মুছে দেয়? আর হজ যা আগে ছিল তা মুছে দেয়?}\lat{”} (মুসলিম ১২১)

\subsubsection*{সংশ্লিষ্ট দলিল}

\begin{itemize}
\item \textbf{সতর্কতা:} \lat{“}ইসলাম যা আগে ছিল তা মুছে দেয়\lat{”} — এটি \textbf{নতুন মুসলিমের} ক্ষেত্রে; \lat{“}\arab{التوبة} \arab{تجب} \arab{ما} \arab{قبلها}\lat{”} (তাওবা যা আগে ছিল তা মুছে দেয়) — \textbf{নবীজির হাদিস হিসেবে প্রমাণিত নয়} (শাইখ সালেহ আল-উসাইমি) — নিচের সতর্কতা টেবিল দেখুন।
\end{itemize}
\medskip\hrule\medskip

\subsection*{হাদিস ৭ — ভুল, ভুলে যাওয়া ও বাধ্য হয়ে করা — ক্ষমা}

\subsubsection*{সূত্র কার্ড}

\begin{table}[h]\centering\small
\begin{tabular}{ll}
বিষয় & বিবরণ \\
\hline
\textbf{বর্ণনাকারী} & ইবনে আব্বাস (রাঃ) \\
\textbf{গ্রন্থ} & সুনানে ইবনে মাজাহ ২০৪৫; সহীহ ইবনে হিব্বান \\
\textbf{অধ্যায়} & কিতাবুত তালাক — \lat{“}বাধ্য ও ভুলে যাওয়া ব্যক্তির তালাক\lat{”} \\
\textbf{সনদের মান} & \textbf{সহীহ} (দারুসসালাম: সহীহ) \\
\hline
\end{tabular}
\end{table}

\subsubsection*{পূর্ণ অনুবাদ}

ইবনে আব্বাস (রাঃ) থেকে বর্ণিত — রাসূলুল্লাহ (সা.) বলেছেন: \textbf{\lat{“}আল্লাহ আমার উম্মতের ভুল, ভুলে যাওয়া এবং যা করতে বাধ্য করা হয়েছে — তা ক্ষমা করে দিয়েছেন।\lat{”}}

(সুনানে ইবনে মাজাহ ২০৪৫)

\subsubsection*{সংশ্লিষ্ট দলিল}

\begin{itemize}
\item \textbf{আয়াত (২:২৮৬):} \lat{“}আল্লাহ কাউকে তার সামর্থ্যের বাইরে বোঝা চাপান না।\lat{”}
\item এই হাদিস প্রমাণ করে — \textbf{কিছু পাপ আল্লাহ মূলতই ক্ষমা করে দিয়েছেন} (ভুল-ভুলে যাওয়া-বাধ্য)। এ ধরনের পাপ কিয়ামতে আর তোলা হবে না।
\end{itemize}
\medskip\hrule\medskip

\subsection*{হাদিস ৮ — তাওবার পরের নেক আমল: পাপ মুছে দেয় (১১:১১৪-এর শানে নুজুল)}

\subsubsection*{সূত্র কার্ড}

\begin{table}[h]\centering\small
\begin{tabular}{ll}
বিষয় & বিবরণ \\
\hline
\textbf{বর্ণনাকারী} & আবদুল্লাহ ইবনে মাসউদ (রাঃ) \\
\textbf{গ্রন্থ} & সহীহ বুখারি ৪৬৮৭; সহীহ মুসলিম ২৭৬৩ \\
\textbf{অধ্যায়} & বুখারি: কিতাবুত তাফসীর — সূরা হুদ \\
\textbf{সনদের মান} & \textbf{সহীহ} (বুখারি-মুসলিম) \\
\hline
\end{tabular}
\end{table}

\subsubsection*{পূর্ণ অনুবাদ}

ইবনে মাসউদ (রাঃ) বলেন: এক ব্যক্তি একজন নারীকে চুমু দিয়েছিল (যিনা করেনি)। সে নবীজির (সা.) কাছে এসে বিষয়টি জানাল। তখন \lat{“}দিনের দুই প্রান্তে নামাজ কায়েম করো... ভালো কাজ মন্দ কাজ মুছে দেয়\lat{”} (১১:১১৪) নাজিল হলো। লোকটি জিজ্ঞেস করল: \lat{“}হে আল্লাহর রাসূল! এটি কি শুধু আমার জন্য?\lat{”} নবীজি (সা.) বললেন: \textbf{\lat{“}বরং আমার উম্মতের মধ্যে যে এ অনুযায়ী আমল করবে — তার জন্যই এটি।\lat{”}}

(সহীহ বুখারি ৪৬৮৭; সহীহ মুসলিম ২৭৬৩)

\subsubsection*{সংশ্লিষ্ট দলিল}

\begin{itemize}
\item \textbf{আয়াত (১১:১১৪):} \lat{“}নিশ্চয়ই ভালো কাজ মন্দ কাজ মুছে দেয়।\lat{”}
\item \textbf{হাদিস (বুখারি ৬৫৩৬; মুসলিম ২৮৭৬):} সহজ হিসাব = আরয।
\end{itemize}
\medskip\hrule\medskip

\subsection*{সতর্কতা টেবিল — প্রচলিত কিন্তু দুর্বল/অসমর্থিত কনটেন্ট}

\begin{table}[h]\centering\small
\begin{tabular}{llll}
\# & প্রচলিত বিষয় & আসল অবস্থা & ব্যাখ্যা ও উৎস \\
\hline
১ & \textbf{\lat{“}\arab{التوبة} \arab{تجب} \arab{ما} \arab{قبلها}\lat{”} — তাওবা যা আগে ছিল তা মুছে দেয়} — হাদিস হিসেবে প্রচার & \textbf{হাদিস নয়} & শাইখ সালেহ আল-উসাইমি: এটি নবীজির হাদিস নয়, নবীজি থেকে বর্ণিত শুধু \lat{“}ইসলাম/হিজরত/হজ যা আগে ছিল তা মুছে দেয়\lat{”} (মুসলিম ১২১)। অর্থটি কুরআন (২৫:৭০) ও সহীহ হাদিসে সমর্থিত — তবে বাক্যটি হাদিস নয়। \\
২ & \textbf{\lat{“}\arab{التائب} \arab{من} \arab{الذنب} \arab{كمن} \arab{لا} \arab{ذنب} \arab{له}\lat{”} — তাওবাকারী পাপহীনের মতো} & \textbf{দুর্বল (মতভেদ)} & ইবনে মাজাহ ৪২৫০ — দারুসসালাম: দুর্বল; আলবানি (সহীহ ইবনে মাজাহ)-এ সহীহ বলেছেন। মতভেদ থাকায় হাদিস হিসেবে দৃঢ়ভাবে বলা যাবে না; তবে অর্থ ২৫:৭০ ও নাজওয়া হাদিসে সমর্থিত। \\
৩ & \textbf{তিরমিযি ২৪৫৯ — \lat{“}কাইস (চতুর) সে, যে নিজের নফসকে হিসাব নেয়\lat{”}} & \textbf{দুর্বল} & দারুসসালাম: দুর্বল (আবু বকর ইবনে আবি মারইয়ামের কারণে)। উমার (রাঃ)-এর \lat{“}তোমরা নিজেদের হিসাব নাও নিজেদের হিসাব নেওয়ার আগে\lat{”} — তিরমিযিতে বর্ণিত উক্তি, কিন্তু সহীহ সনদে প্রমাণিত নয়। শিক্ষামূলক, হাদিসের মানে নয়। \\
৪ & \textbf{দীর্ঘ নাজওয়া সংস্করণ — \lat{“}বান্দার আমলনামা থেকে পাপ মুছে যেতে দেখে, ফেরেশতারাও জানে না\lat{”}} & \textbf{বুখারিতে নেই — সতর্কতা} & প্রচলিত দীর্ঘ সংস্করণ (\lat{hadithanswers} যাচাই করেছে) — বুখারিতে এমন বর্ণনা পাওয়া যায় না। সহীহ সংস্করণ: বুখারি ২৪৪১ / মুসলিম ২৭৬৮ (উপরে দেওয়া)। মূল সহীহ বর্ণনা দিয়েই কাজ চলবে। \\
৫ & \textbf{\lat{“}ক্ষমা করা পাপ কিয়ামতে সবার সামনে প্রকাশ হবে\lat{”}} & \textbf{ভুল ধারণা} & নাজওয়া হাদিস (বুখারি ২৪৪১) ও ইবনে হাজারের ব্যাখ্যা অনুযায়ী মুমিনের পাপ গোপন রাখা হয়; প্রকাশ হয় কাফির-মুনাফিক ও বিশ্বাসঘাতকদের (বুখারি ৩১৮৮)। ৮৬:৯-এর \lat{“}গোপন প্রকাশ\lat{”} মানে আল্লাহর কাছে প্রকাশ — সবার সামনে নয়। \\
৬ & \textbf{\lat{“}তাওবা করলে মানুষের হকও মাফ হয়ে যায়\lat{”}} & \textbf{ভুল} & মুফলিস হাদিস (মুসলিম ২৫৮১) — মানুষের হক ফেরত/ক্ষমা আদায় ছাড়া ক্ষমা হয় না। তাওবা আল্লাহর হক ক্ষমা করে; মানুষের হক ভিন্ন বিষয়। \\
\hline
\end{tabular}
\end{table}

\medskip\hrule\medskip

\subsection*{মূল সিদ্ধান্ত (দলিলের আলোকে)}

\begin{enumerate}
\item \textbf{ক্ষমা করা পাপ আর হিসাবের বোঝা নয়} — নাজওয়া হাদিসে ক্ষমার পর তাকে \lat{“}নেকির আমলনামা\lat{”} দেওয়া হয় (বুখারি ২৪৪১)।
\item \textbf{সহজ হিসাব = আরয (পেশকরণ)} — পুঙ্খানুপুঙ্খ জিজ্ঞাসাবাদ নয়; পুঙ্খানুপুঙ্খ জিজ্ঞাসাবাদ মানেই শাস্তি (বুখারি ৬৫৩৬)।
\item \textbf{মুমিনের পাপ গোপন থাকে} — দুনিয়ায় গোপন রাখলে আখিরাতেও গোপন (বুখারি ২৪৪১; মুসলিম ২৬৯৯); পাপ প্রকাশকারী (মুজাহির) ক্ষমার সাধারণ বিধানের বাইরে (বুখারি ৬০৬৯)।
\item \textbf{মানুষের হক ভিন্ন} — ফেরত/ক্ষমা আদায় ছাড়া ক্ষমা নয় (মুসলিম ২৫৮১); বিশ্বাসঘাতকের পতাকা (বুখারি ৩১৮৮)।
\item \textbf{নতুন মুসলিমের পূর্ব পাপ ইসলামে মুছে যায়} (মুসলিম ১২১); ভুল-ভুলে যাওয়া-বাধ্য হয়ে করা পাপ মূলত ক্ষমা (ইবনে মাজাহ ২০৪৫)।
\end{enumerate}

\end{document}
